Congestion Control Algorithms are critical for regulating data flow and optimizing resource utilization in modern networks. 
This paper provides a comprehensive overview of contemporary Congestion Control Algorithms, tracing the evolution from foundational protocols to highly specialized modern advancements. 
We first establish a theoretical framework by examining traditional loss-based and delay-based algorithms, including TCP Cubic, TCP Vegas, and FAST TCP, while detailing the metrics used to evaluate their performance and fairness.  

The core of this work surveys recent state-of-the-art developments, specifically focusing on BBR and Copa, alongside their respective extensions—AQDC-BBR, BBR-ES, Copa+, and Copa-D—which aim to mitigate issues like "bufferbloat" and RTT unfairness. 
We further explore emerging architectural paradigms such as Performance-oriented Congestion Control (PCC), which utilizes online learning, PowerTCP, which leverages In-band Network Telemetry (INT) for datacenters, and C4, designed for real-time media over unstable wireless links. 
Finally, we provide a comparative analysis of these algorithms based on their design objectives and performance data. 
Our synthesis identifies a trend toward high specialization for distinct environments—such as 4G and 5G mobile networks and datacenters—while highlighting the ongoing challenge of maintaining architectural simplicity and fairness in an increasingly diverse multi-protocol global internet.
